\section{乱序核心中的异常优先级与恢复入口}

乱序执行允许多个功能单元在同一拍报告事件：较早 load 发现缺页，较晚除法发现除零，另一条分支同时发现误预测，外部中断也可能到达。控制器不能按“谁先拉高信号”决定处理顺序，而要先把同步异常记录到各自 ROB 项，再按程序顺序选择最早不能提交的操作。误预测只清除分支之后的错误路径；异常要保留异常之前已经提交的状态；外部中断通常等待一个体系结构规定的可中断边界。

\begin{longtable}{L{0.20\linewidth}L{0.25\linewidth}L{0.27\linewidth}L{0.18\linewidth}}
\toprule
事件 & 首次发现位置 & 进入控制面的条件 & 恢复边界 \\
\midrule
分支误预测 & 分支执行单元 & 分支结果与预测不符 & 分支之后的年轻操作 \\
缺页或权限异常 & 地址转换/访存单元 & 故障 load/store 到达 ROB 队首 & 故障指令本身不提交 \\
算术同步异常 & 对应执行单元 & 异常项到达 ROB 队首 & 异常指令本身不提交 \\
外部中断 & 中断控制器与提交端 & 当前优先级允许且到达可中断边界 & 已退休指令之后 \\
机器检查 & Cache、内存或平台 RAS & 依据错误可恢复性和体系结构规则 & 可能是当前指令、核心或整机 \\
\bottomrule
\end{longtable}

恢复动作要有唯一所有者。提交端冻结新的退休，选定最早事件，保存返回 PC、原因和补充状态，再向前端、重命名、发射、LSQ 和执行单元广播清除范围。各结构用 ROB 序号判断条目是否比恢复点年轻；不能只依赖“清空整个队列”，因为较早且仍有效的操作可能需要保留。恢复完成的边界是所有相关队列的代际或尾指针已经切换，新的取指才可以沿异常向量或正确分支目标继续。

\topic{加载重放与流水线冲刷不是同一种恢复}

load 可能因为多种原因需要再次执行：更早 store 的地址后来证明重叠；Cache line 在访问期间被一致性探测失效；地址翻译或权限状态发生变化；所需端口被占用。若错误只影响一条 load 及其数据依赖链，核心可以做选择性重放；若状态无法局部证明，例如页表权限异常或分支路径错误，则要从明确的程序顺序边界冲刷更年轻操作。

\begin{longtable}{L{0.20\linewidth}L{0.27\linewidth}L{0.27\linewidth}L{0.16\linewidth}}
\toprule
恢复方式 & 保留什么 & 重新执行什么 & 典型代价 \\
\midrule
端口重试 & 所有队列与映射 & 尚未获得端口的操作 & 若干拍排队 \\
选择性重放 & 无关操作和架构映射 & 违规 load 及其依赖者 & 依赖链长度 \\
分支冲刷 & 分支及更早操作 & 正确路径后续指令 & 前端重填时间 \\
异常恢复 & 异常之前已提交状态 & 修复后从故障指令或规定返回点开始 & 软件处理与流水线重建 \\
\bottomrule
\end{longtable}

选择性重放需要跟踪消费者关系。若 L1 的结果已经唤醒 A2，A2 又唤醒 A3，那么只重做 L1 会让 A2、A3 保留基于旧值的结果。实现可以保存依赖图、给结果附加代际，或者保守地清除从 L1 之后的一段窗口。恢复粒度越精细，正常吞吐越高，但状态与验证复杂度也越大。

\topic{提交宽度、ROB 容量与隐藏时延}

ROB 容量决定核心最多能跨过多长的停顿继续寻找独立工作。若一次 Cache miss 需要 200 拍，后续平均每拍进入 4 条微操作，而 ROB 只有 192 项，那么不到 48 拍窗口就会装满；即使 miss 后面还有大量独立指令，前端也会因 ROB 无空位而停止。扩大 ROB 可以看到更远的独立工作，却会增加重命名、旁路、唤醒、恢复和功耗成本。

提交宽度是另一个独立上限。假设执行端平均每拍完成 6 条，但提交端最多退休 4 条，ROB 长期净增加 2 项/拍，最终仍会满。短时间的完成峰值可以由 ROB 吸收，长期吞吐必须满足：
\[
\text{IPC} \leq \min(\text{取指宽度},\text{译码宽度},\text{重命名宽度},
\text{执行端持续吞吐},\text{提交宽度}).
\]

ROB 还会受到队首阻塞。较早的长时延除法或 miss 未完成时，后面数十条短操作即使全部执行完成也不能越过它退休；这种现象称为 head-of-line blocking。性能计数器若只报告“执行单元空闲”容易误判，应同时观察 ROB 占用、提交停顿原因、前端停顿和未决 miss 数。

\topic{内存顺序违规的性能账}

等待所有更早 store 地址确定后再执行 load 最安全，却会丢失大量并行性；总是假设无依赖并立即执行最激进，却可能频繁重放。现代核心通常使用 memory dependence predictor，根据 load 的历史判断它是否应等待某个更早 store。预测器的收益可以用期望代价衡量。

设等待策略每次固定损失 6 拍；推测策略在无冲突时不损失，在冲突时以 40 拍完成检测、清除与重放。若冲突概率为 $p$，推测的期望代价约为 $40p$。当 $40p<6$，即 $p<15\%$ 时，推测平均更优；但这只是平均值，安全性仍依赖每次违规都能被检测并正确恢复。

同一计算也说明为什么“重放次数少”不等于代价小。若违规 load 位于长依赖链根部，单次重放会清除大量工作；若位于窗口末端，影响可能只有几条。性能报告应同时记录违规次数、每次清除的平均操作数、恢复拍数和对提交端造成的空洞。

\section*{章末练习}

\begin{chapterexercise}{事件优先级}
同一拍中，ROB 序号 40 的分支发现误预测，序号 35 的 load 报缺页，序号 52 的除法报除零。说明提交控制器最终选择哪个入口，以及其余两个事件如何处理。
\exercisecheck{学习目标}{同步事件按程序顺序仲裁，较年轻事件可能随冲刷消失}
\end{chapterexercise}

\begin{chapterexercise}{选择性重放}
L1 产生 A2 的源，A2 产生 A3 的源；另有与三者无关的 M4。L1 因内存顺序违规需要重放。最小正确重放集合是什么？为什么不能只重放 L1？
\exercisecheck{学习目标}{恢复范围必须覆盖错误值的传递闭包}
\end{chapterexercise}

\begin{chapterexercise}{窗口容量}
一次 miss 为 180 拍，前端平均每拍向 ROB 放入 3 条微操作，ROB 在 miss 发生时已有 40 项且总容量 256 项。忽略退休，多少拍后 ROB 会满？它能否覆盖完整 miss？
\exercisecheck{学习目标}{窗口项数与可隐藏拍数之间的换算}
\end{chapterexercise}

\begin{chapterexercise}{提交瓶颈}
核心平均每拍完成 5 条操作，提交宽度为 3，ROB 初始空闲 120 项。若完成率持续不变，约多少拍后因为提交不足填满这些空位？
\exercisecheck{学习目标}{区分瞬时完成能力与长期退休吞吐}
\end{chapterexercise}

\begin{chapterexercise}{依赖预测}
等待未知 store 地址固定损失 8 拍；错误推测一次代价 50 拍。求推测优于总等待的冲突概率上限，并说明该比较为什么不能代替正确性验证。
\exercisecheck{学习目标}{用期望代价选择策略，同时保留违规检测}
\end{chapterexercise}

\section*{参考解答与要点}

\begin{chapteranswer}{事件优先级}
序号 35 的缺页是最早不能提交的同步事件，提交端应在它到达队首时进入缺页异常。序号 40 的分支和序号 52 的除法都更年轻，会随异常恢复被清除；返回后重新取指时，它们是否再次发生由新的实际执行决定。
\answercheck{必须掌握的要点}{按 ROB 顺序选择最早事件，不按报告信号到达先后选择}
\end{chapteranswer}

\begin{chapteranswer}{选择性重放}
最小集合是 L1、A2、A3；M4 与错误值无依赖，可以保留。只重放 L1 会留下 A2、A3 基于旧值计算的结果，使错误继续传播。
\answercheck{必须掌握的要点}{依赖者要递归失效，无关操作无需清除}
\end{chapteranswer}

\begin{chapteranswer}{窗口容量}
剩余空位为 $256-40=216$ 项，以 3 项/拍进入，$216/3=72$ 拍后填满。72 拍远小于 180 拍，因此该 ROB 不能单独覆盖完整 miss，除非期间仍有较早操作退休释放空间。
\answercheck{必须掌握的要点}{窗口覆盖时间由剩余项数除以进入速率得到}
\end{chapteranswer}

\begin{chapteranswer}{提交瓶颈}
ROB 每拍净增加 $5-3=2$ 项，120 个空位约在 $120/2=60$ 拍后耗尽。之后完成端必须因 ROB 或后续队列背压降低速率。
\answercheck{必须掌握的要点}{长期完成率不能持续高于提交率}
\end{chapteranswer}

\begin{chapteranswer}{依赖预测}
推测期望代价为 $50p$，优于 8 拍等待要求 $50p<8$，即 $p<16\%$。无论概率多低，硬件仍必须比较地址、发现违规并重放；期望代价只决定是否值得推测，不证明推测结果正确。
\answercheck{必须掌握的要点}{性能预测器可以猜，正确性检查不能省略}
\end{chapteranswer}
